\chapter{Sauts et modèles de Lévy}\label{ch:dv:jumps-and-levy-models}

Une société de biotechnologie publiera jeudi les résultats de son essai clinique principal. Ses options à la monnaie d'une semaine
se négocient à une volatilité implicite de 160~\%. Personne ne croit que le titre diffusera à
160~\% par an. Il évoluera au rythme habituel de 35~\% jusqu'à jeudi, puis fera un saut à la hausse d'environ 22~\% ou à la baisse d'environ 15~\%,
avant de reprendre sa dérive. Une diffusion a besoin d'une volatilité énorme pour atteindre le même prix, et elle valorise alors mal
tous les autres strikes et toutes les autres échéances, et donne une mauvaise couverture. Un modèle à sauts n'a besoin que du seul
paramètre sur lequel le desk a réellement une opinion, à savoir l'ampleur du mouvement. Ce chapitre présente les éléments
qui montrent que les options sur indices et sur actions individuelles contiennent des sauts, ainsi que les trois modèles de diffusion à sauts utilisés en pratique, ceux de Merton, de Kou et
de Bates. Il aborde ensuite les modèles de Lévy purs en sauts, ce qu'aucune couverture sur le sous-jacent ne peut éliminer, et les
situations dans lesquelles les sauts modifient le prix d'une transaction.

\section{Pourquoi les diffusions échouent aux échéances courtes}

Dans une diffusion, quelle que soit sa volatilité, le prix évolue de façon continue. Sur un court intervalle de temps $T$, son log-rendement est
approximativement normal, d'écart type $\sigma\sqrt T$, de sorte qu'une option dont le strike est fixé à une distance donnée hors de la monnaie
a besoin d'un mouvement de plusieurs écarts types pour être exercée. Son prix décroît plus vite que toute puissance de $T$ lorsque
l'échéance se réduit~: la queue gaussienne $e^{-k^2/(2\sigma^2T)}$ décroît exponentiellement en $1/T$. Avec des sauts, un seul saut peut faire passer le prix
au-delà du strike à tout instant. La probabilité d'un saut avant $T$ vaut environ
$\lambda T$ pour une intensité de saut $\lambda$, de sorte que le prix de l'option soit proportionnel à $T$ pour $T$ petit~:
\[
P_{\mathrm{OTM}}(T)\approx\lambda T\,\E\bigl[(K-S_0e^{J})^+\bigr]\quad\text{pour une option de vente, }T\to0,
\]
avec $J$ la taille logarithmique d'un saut. À la monnaie, les deux composantes contribuent. La diffusion donne un prix
proportionnel à $\sqrt T$, qui domine. Carr et Wu ont transformé la différence entre ces vitesses en un test
de la nature du processus qui sous-tend les options, et ont trouvé sur les options sur S\&P~500 à la fois une composante continue et une
composante de sauts.

\begin{example}[Un put à 10~\% lorsque l'échéance se réduit]\label{ex:dv:jumps-and-levy-models:carrwu}
Prenons le modèle de Merton ajusté ci-dessous à la surface du chapitre~9 ($\sigma=9.8\%$, $\lambda=3.16$ sauts par an de moyenne
de taille logarithmique $-5.8\%$). Son put hors de la monnaie à 10~\%, divisé par son échéance, tend vers
$\lambda\E[(K-S_0e^J)^+]=1.00$ par an~: à un jour, le rapport vaut 1.02. Sous Black--Scholes avec la même
variance totale, 16.4~\% par an, le même put vaut $3\times10^{-36}$ à un jour et $6\times10^{-9}$ à cinq jours
(la \cref{fig:dv:jumps-and-levy-models:short}, à gauche).
\end{example}

Le même argument explique le smile de volatilité. La volatilité implicite à courte échéance d'une diffusion est plate, car
sa distribution à courte échéance est normale. Celle d'un saut est pentue, car sa distribution à courte échéance est un mélange
de «~rien ne s'est produit~» et «~un saut s'est produit~», avec des queues épaisses. Pour un processus de Lévy, la classe qui contient
tous les modèles de ce chapitre sauf celui de Bates, les cumulants du log-rendement croissent linéairement en $T$. Dès lors,
l'asymétrie varie comme $T^{-1/2}$ et l'excès de kurtosis comme $T^{-1}$. Le smile est le plus pentu aux courtes
échéances et s'estompe à mesure que le théorème central limite prend le relais.

\begin{omfigure}
\begin{tikzpicture}
\begin{groupplot}[group style={group size=2 by 1,horizontal sep=1.6cm},
  width=6.6cm,height=5cm,
  tick label style={font=\small},label style={font=\small},grid=major,grid style={gray!25},table/col sep=comma]
\nextgroupplot[xmode=log,ymode=log,xlabel={échéance $T$ (années)},ylabel={prix du put de strike 90},
  xmin=0.0025,xmax=0.28,ymin=1e-8,ymax=3,
  title={\small put à 10~\%},
  legend columns=1,legend cell align=left,legend style={font=\footnotesize,at={(0.5,-0.32)},anchor=north,draw=none,fill=none}]
\addplot[omThm,very thick] table[x=t,y=jump]{figdata/derivatives/13-jumps-and-levy-models/otm_decay.csv};
\addplot[omDef,very thick,dashed,unbounded coords=jump] table[x=t,y=bs]{figdata/derivatives/13-jumps-and-levy-models/otm_decay.csv};
\addplot[gray,thick,dotted] table[x=t,y=limit]{figdata/derivatives/13-jumps-and-levy-models/otm_decay.csv};
\legend{Merton,{Black--Scholes, même variance},{$\lambda T\,\E[(K-S_0e^J)^+]$}}
\nextgroupplot[xmode=log,ymode=log,xlabel={échéance $T$ (années)},ylabel={skew à la monnaie $|\psi(T)|$},
  xmin=0.015,xmax=2.6,ymin=0.025,ymax=5,
  xtick={0.02,0.1,0.5,2},xticklabels={0.02,0.1,0.5,2},ytick={0.03,0.1,0.3,1,3},yticklabels={0.03,0.1,0.3,1,3},
  title={\small structure par terme du skew},
  legend columns=2,legend cell align=left,legend style={font=\footnotesize,at={(0.5,-0.32)},anchor=north,draw=none,fill=none}]
\addplot[black,only marks,mark=*,mark size=1.8pt] table[x=t,y=market]{figdata/derivatives/13-jumps-and-levy-models/skews.csv};
\addplot[omThm,very thick] table[x=t,y=merton]{figdata/derivatives/13-jumps-and-levy-models/skews.csv};
\addplot[omDef,thick,dashed] table[x=t,y=vg]{figdata/derivatives/13-jumps-and-levy-models/skews.csv};
\addplot[omProp,thick] table[x=t,y=heston]{figdata/derivatives/13-jumps-and-levy-models/skews.csv};
\addplot[gray,very thick,dotted] table[x=t,y=bates]{figdata/derivatives/13-jumps-and-levy-models/skews.csv};
\legend{marché,Merton,variance gamma,Heston,Bates}
\end{groupplot}
\end{tikzpicture}

\omcaption{Les échéances courtes séparent les sauts de la diffusion. À gauche~: un put hors de la monnaie à 10~\% sous le
modèle de Merton ajusté est proportionnel à son échéance~; sous Black--Scholes avec la même variance, il s'annule
plus vite que toute puissance. À droite~: le skew à la monnaie de la surface du chapitre~9, comparé à Merton et variance gamma
ajustés à un mois, Heston ajusté sur la surface (chapitre~10) et Bates ajusté aux huit échéances. Données~:
le tutoriel.}\label{fig:dv:jumps-and-levy-models:short}
\end{omfigure}

\section{Diffusions à sauts~: Merton, Kou, Bates}

Une diffusion à sauts ajoute au mouvement brownien du chapitre~3 un processus de Poisson composé (le livre~4,
chapitre~6)~: les sauts surviennent au taux $\lambda$ et chacun multiplie le prix par $e^{J}$, les $J$ étant indépendants et
identiquement distribués. Pour que le prix actualisé reste une martingale, la dérive est abaissée de
l'espérance du saut relatif, $\lambda\bar k$ avec $\bar k=\E[e^J]-1$. Avec $x_T=\ln(S_T/F_T)$ et $\phi_J$ la
fonction caractéristique de $J$, tout modèle de ce type vérifie
\[
\E\bigl[e^{iux_T}\bigr]=\exp\Bigl(T\bigl(-\tfrac12\sigma^2(iu+u^2)-iu\lambda\bar k+\lambda(\phi_J(u)-1)\bigr)\Bigr),
\]
la forme de Lévy--Khintchine (livre~4, chapitre~6) avec une partie gaussienne et une mesure de sauts finie. Les trois
modèles diffèrent par la loi de $J$ et par le fait que la volatilité soit elle-même aléatoire.

\begin{definition}[Modèle de diffusion à sauts de Merton]\label{def:dv:jumps-and-levy-models:merton}
Le \emph{modèle de diffusion à sauts de Merton}\index{modèle de diffusion à sauts de Merton} est la diffusion à sauts dont les log-sauts sont normaux,
de loi $J\sim\mathcal N(\mu_J,\delta^2)$, donc $\phi_J(u)=e^{iu\mu_J-\frac12\delta^2u^2}$ et
$\bar k=e^{\mu_J+\frac12\delta^2}-1$. Ses paramètres sont $(\sigma,\lambda,\mu_J,\delta)$.
\end{definition}

Le modèle de Merton admet un prix sous forme fermée. Conditionnons sur le nombre $n$ de sauts avant $T$, qui suit une loi de Poisson de
moyenne $\lambda T$. Sachant $n$, le log-rendement est normal, de sorte que l'option est un prix de Black avec une forward décalée et
une variance plus grande~:
\[
C=\sum_{n\ge0}e^{-\lambda T}\frac{(\lambda T)^n}{n!}\,C_{\mathrm{Black}}\Bigl(F\,e^{n(\mu_J+\frac12\delta^2)-\lambda\bar kT},\,K,\,T,\,
\sqrt{\sigma^2+n\delta^2/T}\Bigr).
\]
Le tutoriel le vérifie avec le prix de Fourier de l'intégrale du chapitre~10, à quelques millionièmes près.

\begin{definition}[Modèle de Kou]\label{def:dv:jumps-and-levy-models:kou}
Le \emph{modèle de Kou}\index{modèle de Kou} est la diffusion à sauts dont les log-sauts suivent une loi double exponentielle~: avec probabilité
$p$ le saut est à la hausse et exponentiel de taux $\eta_1>1$, sinon il est à la baisse et exponentiel de taux $\eta_2$, donc
$\phi_J(u)=p\,\frac{\eta_1}{\eta_1-iu}+(1-p)\frac{\eta_2}{\eta_2+iu}$. Ses paramètres sont
$(\sigma,\lambda,p,\eta_1,\eta_2)$.
\end{definition}

Kou a retenu la double exponentielle pour deux raisons. Elle donne des queues asymétriques, plus lourdes que la loi normale de part et d'autre,
et la loi exponentielle est sans mémoire~: le dépassement d'une barrière par un saut suit à nouveau une loi exponentielle,
ce qui donne des prix analytiques aussi bien pour les options dépendantes de la trajectoire que pour les vanilles. La condition $\eta_1>1$ garantit que $\E[e^J]$ est fini.

\begin{definition}[Modèle de Bates]\label{def:dv:jumps-and-levy-models:bates}
Le \emph{modèle de Bates}\index{modèle de Bates} est le \omterm{def:dv:stochastic-volatility:sv}{modèle à volatilité stochastique} de Heston (chapitre 10) auquel on ajoute des sauts de Merton
dans le prix. La fonction caractéristique est le produit de celle de Heston et du facteur de saut, et les
paramètres sont $(v_0,\kappa,\bar v,\eta,\rho,\lambda,\mu_J,\delta)$.
\end{definition}

La répartition des rôles est le point essentiel du \omterm{def:dv:jumps-and-levy-models:bates}{modèle de Bates}. Les sauts produisent le skew court terme. La volatilité
stochastique produit le skew long terme et la structure par terme. Aucun des deux ne peut faire le travail de l'autre~: le skew dû aux sauts s'estompe
en $T^{-1/2}$, et le skew court terme de Heston se stabilise (chapitre 10).

\begin{example}[Trois modèles sur une échéance, un seul sur toutes]\label{ex:dv:jumps-and-levy-models:fits}
Ajusté au smile à un mois de la surface du chapitre 9 (neuf strikes à moins de deux écarts types), Merton donne
$\sigma=9.8\%$, $\lambda=3.16$, $\mu_J=-5.8\%$, $\delta=4.6\%$ avec une erreur quadratique moyenne de 0.18 point
de volatilité. Kou donne $\sigma=8.4\%$, $\lambda=11.7$, $p=0.24$, $\eta_1=85$, $\eta_2=30.4$, avec une erreur de 0.06.
La variance gamma (section suivante) présente une erreur de 0.24 (la \cref{fig:dv:jumps-and-levy-models:fits}). Reportés sur
les autres échéances, les trois modèles échouent. À un an, le skew de Merton vaut $-0.07$ contre $-0.25$ sur le marché. À une semaine, celui de la variance
gamma vaut $-3.66$ contre $-1.44$ (la \cref{fig:dv:jumps-and-levy-models:short}). Bates, ajusté aux huit échéances
allant d'une semaine à deux ans, présente une erreur de 0.34 point, avec $v_0=0.018$, $\kappa=2.06$, $\bar v=0.062$,
$\eta=0.85$, $\rho=-0.70$ et sauts $\lambda=2.34$, $\mu_J=-3.6\%$, $\delta=3.0\%$. Son skew vaut $-1.36$ à une semaine
et $-0.23$ à un an.
\end{example}

\begin{omfigure}
\begin{tikzpicture}
\begin{axis}[width=10.5cm,height=5cm,
  xlabel={strike},ylabel={volatilité implicite à un mois (\%)},
  tick label style={font=\small},label style={font=\small},
  xmin=87,xmax=114,ymin=10,ymax=24,grid=major,grid style={gray!25},
  legend columns=3,legend cell align=left,legend style={font=\footnotesize,at={(0.5,-0.3)},anchor=north,draw=none,fill=none,/tikz/every even column/.append style={column sep=8pt}},
  table/col sep=comma]
\addplot[black,only marks,mark=*,mark size=1.8pt] table[x=k,y=market]{figdata/derivatives/13-jumps-and-levy-models/fit_market.csv};
\addplot[omThm,very thick] table[x=k,y=merton]{figdata/derivatives/13-jumps-and-levy-models/fit_models.csv};
\addplot[omDef,thick,dashed] table[x=k,y=kou]{figdata/derivatives/13-jumps-and-levy-models/fit_models.csv};
\addplot[omProp,thick,dotted] table[x=k,y=vg]{figdata/derivatives/13-jumps-and-levy-models/fit_models.csv};
\addplot[gray,thick] table[x=k,y=heston]{figdata/derivatives/13-jumps-and-levy-models/fit_models.csv};
\legend{cotations du marché,Merton,Kou,variance gamma,Heston (chapitre 10)}
\end{axis}
\end{tikzpicture}

\omcaption{Smiles à un mois de la surface du chapitre 9~: les neuf cotations ajustées, les trois modèles à sauts ajustés à celles-ci,
et Heston calibré sur toute la surface. Sur une seule échéance, tous les modèles à sauts s'ajustent~; les écarts apparaissent
sur les autres échéances. Données~: le tutoriel.}\label{fig:dv:jumps-and-levy-models:fits}
\end{omfigure}

\section{Modèles de Lévy purs}

Un processus de Lévy comportant une infinité de petits sauts dans chaque intervalle peut se passer de la partie brownienne. Les petits
sauts jouent le rôle de la diffusion, et les grands sauts forment les queues. Le modèle le plus connu de ce type modifie
l'horloge d'un mouvement brownien.

\begin{definition}[Modèle variance gamma]\label{def:dv:jumps-and-levy-models:vg}
Le \emph{modèle variance gamma}\index{modèle variance gamma} fait courir un mouvement brownien avec dérive $\theta$ et
volatilité $\sigma$ sur une horloge aléatoire $g_t$, un processus gamma de taux moyen un et de taux de variance $\nu$ (un
subordinateur, livre~4, chapitre~6)~: $x_t=\omega t+\theta g_t+\sigma W_{g_t}$, avec $\omega$ choisi de sorte que
$\E[e^{x_t}]=1$. Sa fonction caractéristique est
$\E[e^{iux_t}]=e^{iu\omega t}\bigl(1-iu\theta\nu+\frac12\sigma^2\nu u^2\bigr)^{-t/\nu}$, et
$\omega=\frac1\nu\ln\bigl(1-\theta\nu-\frac12\sigma^2\nu\bigr)$.
\end{definition}

Les trois paramètres ont un sens direct. $\sigma$ fixe l'échelle, $\nu$ l'aplatissement (la variabilité
de l'horloge), et $\theta$ l'asymétrie (la dérive que porte l'horloge). Madan, Carr et Chang ont constaté que, sur les données du
S\&P~500, la densité statistique était presque symétrique alors que la \omterm{def:dv:implied-volatility-and-its-surface:density}{densité risque-neutre} était négativement asymétrique
et plus aplatie. Le skew des options sur indices, dans cette lecture, est davantage un prix du risque de baisse que
une caractéristique des rendements réalisés. Le conditionnement sur l'horloge donne une forme fermée. Sachant $g_T=g$ le log-rendement est
normal, donc le call est un prix de Black moyenné sur la loi gamma de $g$, que la construction calcule par quadrature
et vérifie par rapport au prix de Fourier. Le modèle CGMY de Carr, Geman, Madan et Yor généralise la variance gamma
par un quatrième paramètre qui contrôle l'activité des petits sauts. Il couvre les cas d'activité finie et infinie,
ainsi que les cas de variation finie et infinie.

\begin{example}[Moments qui s'estompent]\label{ex:dv:jumps-and-levy-models:moments}
Pour le modèle de Merton ajusté ci-dessus, l'asymétrie du log-rendement vaut $-2.90$ à une semaine, $-1.39$ à un mois
et $-0.40$ à un an, et l'aplatissement en excès 15.2, 3.5 et 0.29~: exactement les lois en $T^{-1/2}$ et $T^{-1}$. Le
\omterm{def:dv:jumps-and-levy-models:vg}{modèle variance gamma} ajusté se situe à quelques pour cent de ces mêmes valeurs, puisque les deux ont été ajustés au même
mois. C'est pourquoi un modèle de Lévy ajusté à une seule échéance ne peut pas ajuster les autres. Le skew du marché décroît
en $T^{-0.44}$ (chapitre 12), proche de l'exposant de Lévy $-\frac12$ aux échéances courtes, mais aux échéances longues il
est soutenu par la persistance de la volatilité, que n'a aucun processus de Lévy.
\end{example}

\section{Ce qui ne peut pas être couvert}

Un marché avec des sauts de taille aléatoire est incomplet (chapitre 1)~: l'action seule ne permet pas de répliquer une option.
La couverture en delta fonctionne dans une diffusion parce que, sur un petit intervalle, la valeur de l'option est linéaire
au premier ordre par rapport au mouvement, et le terme d'ordre deux se moyenne grâce au gamma et au \omterm{def:dv:greeks-and-the-hedging-pnl:greeks}{theta}. Un saut n'est pas petit. Sur un
saut, la position couverte gagne ou perd la convexité de l'option sur l'ensemble du mouvement,
$V(Se^J)-V(S)-\Delta S(e^J-1)$, et aucun choix de $\Delta$ ne la supprime pour tous les $J$ à la fois.

\begin{example}[Couverture en delta dans un monde à sauts]\label{ex:dv:jumps-and-levy-models:hedge}
Vendre un call d'un mois à la monnaie au prix de Merton ajusté, soit 1.89 pour un spot de 100, et le couvrir en delta chaque jour
avec le delta de Merton lui-même. Sur 20\,000 trajectoires, le P\&L couvert a une moyenne nulle et un écart type de 0.97 fois
la prime, et une trajectoire sur cent perd plus de 4.1 primes. Dans une diffusion de même variance totale,
couverte de la même manière, l'écart type vaut 0.185 prime et la perte d'une trajectoire sur cent vaut 0.50
(la \cref{fig:dv:jumps-and-levy-models:hedge}, à gauche). La couverture quotidienne est aussi bonne dans les deux mondes entre les sauts~; la
différence tient à la poignée de jours où un saut survient.
\end{example}

L'événement programmé est le cas limite, et un cas révélateur. Si le résultat peut prendre deux tailles connues, l'événement se ramène
à un arbre binomial à une étape (chapitre 2)~: deux états et un instrument de couverture, un \omterm{def:dv:no-arbitrage-and-the-fundamental-theorems:complete}{marché complet}. Le ratio
$\Delta^*=\frac{V_u-V_d}{S_u-S_d}$ supprime entièrement le risque de saut. Ce n'est pas le delta de l'option, qui mesure
la réaction aux petits mouvements. Dès que les tailles sont incertaines, l'arbre comporte plus de branches que d'instruments, et un
résidu subsiste que seules des options peuvent couvrir.

\begin{example}[Couverture à travers le résultat de l'essai]\label{ex:dv:jumps-and-levy-models:event}
Le desk est vendeur du straddle à la monnaie à une semaine sur l'action de biotechnologie, valant 8.86 pour un spot de 50 (une
diffusion à 35~\%, un mouvement à la hausse de 22.1~\% avec la probabilité 0.4, un mouvement à la baisse de 14.8~\%). Il se couvre depuis juste avant
l'annonce jusqu'à juste après. L'écart type du P\&L sur l'événement vaut 1.81, soit 20.4~\% de
la prime, sans couverture. Il vaut 2.01 avec le delta du modèle, $-0.02$, qui ne réagit qu'aux petits mouvements. Il vaut
1.01 (11.4~\%) avec le delta de Black à la volatilité implicite du straddle, $+0.09$. Avec le ratio binomial,
$+0.20$, il est nul. Lorsque la taille de chaque résultat est elle-même incertaine (un écart type de 4~\% en log autour de chacune),
le meilleur ratio unique reste 0.20, mais il laisse 1.45, soit 16.5~\% de la prime
(la \cref{fig:dv:jumps-and-levy-models:hedge}, à droite).
\end{example}

\begin{omfigure}
\begin{tikzpicture}
\begin{groupplot}[group style={group size=2 by 1,horizontal sep=1.6cm},
  width=6.6cm,height=5cm,
  tick label style={font=\small},label style={font=\small},grid=major,grid style={gray!25},table/col sep=comma]
\nextgroupplot[xlabel={P\&L couvert (en primes)},ylabel={part des trajectoires},ymode=log,xmin=-5,xmax=1.5,ymin=1e-4,ymax=1,
  title={\small delta quotidien, un mois},
  legend columns=2,legend style={font=\footnotesize,at={(0.5,-0.32)},anchor=north,draw=none,fill=none}]
\addplot[omThm,very thick,const plot] table[x=x,y=merton]{figdata/derivatives/13-jumps-and-levy-models/hedge_hist.csv};
\addplot[omDef,thick,dashed,const plot] table[x=x,y=diffusion]{figdata/derivatives/13-jumps-and-levy-models/hedge_hist.csv};
\legend{modèle de Merton,diffusion}
\nextgroupplot[xlabel={couverture},ylabel={écart type du P\&L (\% de prime)},ymin=0,ymax=25,xmin=-0.6,xmax=4.6,
  xtick=data,xticklabels={aucune,modèle,Black,binomial,incertain},
  x tick label style={font=\scriptsize,rotate=35,anchor=north east},
  title={\small à travers l'annonce}]
\addplot[omThm,ybar,bar width=12pt,fill=omThm!40] table[x=i,y=sd]{figdata/derivatives/13-jumps-and-levy-models/event_hedge.csv};
\end{groupplot}
\end{tikzpicture}

\omcaption{Ce que laisse une couverture sur l'action. À gauche~: le P\&L d'un call vendu à un mois, couvert chaque jour, dans le modèle de Merton calibré
et dans une diffusion de même variance (échelle logarithmique~: la queue gauche du monde à sauts est la différence).
À droite~: l'écart type du P\&L d'un straddle vendu et couvert à travers une annonce binaire, sans couverture,
le delta du modèle, le delta de Black et le ratio binomial (tailles de résultats connues), et le meilleur ratio lorsque les tailles
sont incertaines~; le ratio binomial le supprime entièrement quand les tailles sont connues. Données~: le
tutoriel.}\label{fig:dv:jumps-and-levy-models:hedge}
\end{omfigure}

\section{Quand les sauts importent}

Les sauts modifient les prix et les couvertures chaque fois que le payoff est sensible aux grands mouvements sur de courtes périodes.
\begin{itemize}
\item \textbf{Événements programmés.} Résultats d'entreprise, résultats d'essais cliniques, élections, décisions des banques centrales~: la \omterm{def:dv:parametrising-the-surface:event}{variance d'événement}
du chapitre 8 est un saut. Un modèle à deux points lit la probabilité et les tailles dans le smile, ce qu'une seule
volatilité implicite ne peut pas faire.
\item \textbf{Options hors de la monnaie à courte échéance.} Leur valeur tient surtout à la probabilité d'un saut
(l'\cref{ex:dv:jumps-and-levy-models:carrwu})~; une diffusion calibrée ailleurs les rend presque gratuites.
\item \textbf{Payoffs sensibles aux gaps.} Une barrière suivie en continu (chapitre 15) peut être franchie par un saut
sans que le cours échange à proximité, de sorte que la couverture prévue par le desk à la barrière n'a jamais lieu. Les digitales et
les autocallables proches de leur déclenchement, et les produits à effet de levier qui doivent être rééquilibrés, portent le même risque de gap.
\item \textbf{Réserves de couverture.} Lorsque la couverture laisse un résidu de saut, le prix doit porter une réserve pour lui
(chapitre 27)~; la taille du résidu dans l'\cref{ex:dv:jumps-and-levy-models:hedge} donne l'ordre de grandeur.
\end{itemize}
Là où ils comptent moins~: les vanilles à longue échéance, dont le smile tient surtout à la dynamique de la volatilité, et les
options à la monnaie, dont le prix tient surtout à la diffusion.

\begin{omfigure}
\begin{tikzpicture}
\begin{groupplot}[group style={group size=2 by 1,horizontal sep=1.6cm},
  width=6.6cm,height=5cm,
  tick label style={font=\small},label style={font=\small},grid=major,grid style={gray!25},table/col sep=comma]
\nextgroupplot[xlabel={cours de l'action à l'échéance},ylabel={densité},xmin=25,xmax=80,ymin=0,
  yticklabel style={/pgf/number format/fixed},
  title={\small densité à une semaine},
  legend columns=1,legend cell align=left,legend style={font=\footnotesize,at={(0.5,-0.32)},anchor=north,draw=none,fill=none}]
\addplot[omThm,very thick] table[x=s,y=event]{figdata/derivatives/13-jumps-and-levy-models/event_density.csv};
\addplot[omDef,thick,dashed] table[x=s,y=lognormal]{figdata/derivatives/13-jumps-and-levy-models/event_density.csv};
\legend{diffusion et événement binaire,{lognormale à 160~\% de l'événement}}
\nextgroupplot[xlabel={strike},ylabel={volatilité implicite à une semaine (\%)},xmin=37,xmax=63,ymin=60,ymax=170,
  title={\small smile à une semaine}]
\addplot[omThm,very thick,mark=*,mark size=1.5pt] table[x=k,y=vol]{figdata/derivatives/13-jumps-and-levy-models/event_smile.csv};
\end{groupplot}
\end{tikzpicture}

\omcaption{Une annonce binaire dans une option à une semaine sur une action à 50~: une diffusion à 35~\% plus un mouvement de
$+22.1\%$ (probabilité 0.4) ou $-14.8\%$. À gauche~: la densité présente deux bosses, qu'aucune lognormale seule
reproduit. À droite~: la volatilité implicite vaut 160~\% à la monnaie et décroît des deux côtés, le «~froncement~» d'une
distribution bimodale. Données~: le code du chapitre.}\label{fig:dv:jumps-and-levy-models:event}
\end{omfigure}

\begin{example}[Lire l'événement dans le smile]\label{ex:dv:jumps-and-levy-models:read}
La volatilité à la monnaie de l'action de biotechnologie, 160.4~\%, lue avec la règle du chapitre 8 face à une diffusion de 35~\%,
donne une \omterm{def:dv:parametrising-the-surface:event}{variance d'événement} de 0.047. L'écart type de la règle pour le mouvement vaut alors 21.7~\%, et son
mouvement absolu espéré vaut 17.3~\%. Le mouvement réel a un écart type de 17.6~\% et une taille absolue espérée de 17.6~\%.
Le straddle évalue le mouvement absolu espéré, que la règle restitue correctement, tandis que l'écart type de la règle
suppose un mouvement normal et surestime le mouvement binaire. Un modèle à deux points calibré sur l'ensemble du smile
(la \cref{fig:dv:jumps-and-levy-models:event}) retrouve la probabilité et les deux tailles. Il le fait même lorsque les
tailles des résultats du marché sont incertaines, comme le montre le problème du week-end.
\end{example}

\section{Tutoriel~: une interface pour les modèles à sauts}\label{tut:dv:jumps-and-levy-models:tut}

\begin{tutorial}
\textbf{Objectif.} Implémenter quatre modèles à sauts et le modèle d'événement derrière une même interface de fonction caractéristique,
vérifier chacun par rapport à une forme fermée, les calibrer sur un smile, et mesurer les résidus de couverture. \textbf{État final~:} les
quatre figures du chapitre et les chiffres du problème du week-end.
\begin{tutsteps}
\item \textbf{Un modèle est sa fonction caractéristique.} Celui de Merton, corrigé pour être une martingale, avec le prix par série
qui le vérifie~:
\omcode{code/firm/jumps/firm_jumps.py}{38}{67}{Modèle de Merton~: fonction caractéristique, cumulants, prix par série.}\label{lst:dv:jumps-and-levy-models:merton}
\item \textbf{Un seul pricer pour tous.} L'intégrale de Lewis du chapitre 10, sur une grille assez longue pour les modèles à sauts purs
aux échéances courtes, dont les fonctions caractéristiques ne décroissent qu'en puissance de $u$~:
\omcode{code/firm/jumps/firm_jumps.py}{175}{182}{Pricer de Fourier utilisant la fonction caractéristique de tout modèle.}\label{lst:dv:jumps-and-levy-models:lewis}
\item \textbf{L'événement programmé} comme modèle, avec un prix exact de mélange à deux points~:
\omcode{code/firm/jumps/firm_jumps.py}{147}{172}{Le modèle d'événement binaire.}\label{lst:dv:jumps-and-levy-models:event}
\item \textbf{Exécuter} \texttt{dv\_jumps.fit\_models()}, \texttt{fit\_bates()}, \texttt{event\_hedge()},
\texttt{hedge\_experiment()} et \texttt{fig\_jumps.py}.
\end{tutsteps}
\textbf{Pour aller plus loin.} Calibrer Merton à une semaine au lieu d'un mois et comparer le skew à un an~; fixer
l'incertitude sur la taille des résultats à 8~\% et relancer la couverture d'événement~; remplacer l'horloge gamma de la variance gamma par une
inverse gaussienne (le modèle normal inverse Gaussian) et comparer les ajustements.
\end{tutorial}

\section{Construction~: modèles à sauts derrière une même interface}\label{bld:dv:jumps-and-levy-models:jumps}

\begin{build}
\bfield{Rôle} Les modèles à sauts de la firme miniature, sous forme de fonctions caractéristiques que le moteur de transformées du
chapitre 24 price et calibre, ainsi que le modèle d'événement programmé qu'utilise le desk d'options autour des annonces.
\bfield{Interface} Dataclasses \texttt{Merton}, \texttt{Kou}, \texttt{Bates}, \texttt{VarianceGamma},
\texttt{EventJump} (avec \texttt{from\_up}), chacune avec \texttt{cf(u, t)} de $\ln(S_t/F_t)$ et, là où elles existent,
\texttt{cumulants(t)}, \texttt{series\_call}, \texttt{conditional\_call}, \texttt{mixture\_call}~;
\texttt{call\_prices(model, fwd, strikes, t)}~; \texttt{implied\_vols}~; \texttt{skew\_kurtosis}~;
\texttt{calibrate(make, z0, quotes, fwd) -> (model, z, rmse)}.
\bfield{Règles} Toute fonction caractéristique vérifie $\phi(0)=1$ et $\phi(-i)=1$ (la condition de martingale),
vérifiée dans les tests~; les paramètres sont calibrés par des transformations qui les gardent admissibles~; la grille de Fourier est
vérifiée par rapport à une forme fermée à l'échéance la plus courte que le desk négocie.
\bfield{Tests d'acceptation} \texttt{code/firm/jumps/tests/}~: la condition de martingale pour chaque modèle à une semaine
et à un an~; Fourier face à la série de Merton, au prix conditionnel de la variance gamma et au mélange d'événement~; sans
sauts donne Black--Scholes, et Bates sans sauts donne Heston~; l'asymétrie et l'aplatissement varient comme $T^{-1/2}$ et
$T^{-1}$~; la calibration retrouve les paramètres connus de Merton.
\bfield{Pour aller plus loin} Les modèles CGMY et normal inverse Gaussian avec la même interface~; les prix de barrière de Kou en forme fermée
servant de contrôle aux moteurs du chapitre 15.
\end{build}

\begin{omsources}
\item R.~C.~Merton, «~Option pricing when underlying stock returns are discontinuous~», \emph{Journal of Financial
Economics} 3(1--2) (1976) 125--144.
\item S.~G.~Kou, «~A jump-diffusion model for option pricing~», \emph{Management Science} 48(8) (2002) 1086--1101.
\item D.~S.~Bates, «~Jumps and stochastic volatility~: exchange rate processes implicit in Deutsche mark options~»,
\emph{Review of Financial Studies} 9(1) (1996) 69--107.
\item D.~B.~Madan, P.~P.~Carr et E.~C.~Chang, «~The variance gamma process and option pricing~», \emph{Review of
Finance} 2(1) (1998) 79--105.
\item P.~Carr, H.~Geman, D.~B.~Madan et M.~Yor, «~The fine structure of asset returns~: an empirical
investigation~», \emph{Journal of Business} 75(2) (2002) 305--333.
\item P.~Carr et L.~Wu, «~What type of process underlies options~? A simple robust test~», \emph{Journal of
Finance} 58(6) (2003) 2581--2610.
\end{omsources}

\section{Exercices}

\begin{exercise}[$\star$]\label{exo:dv:jumps-and-levy-models:1}
Dans le modèle de Merton avec $\lambda=3$ par an, quelle est la probabilité d'au moins un saut en une semaine (une année de 52
semaines)~? Et en un an~?
\end{exercise}

\begin{exercise}[$\star$]\label{exo:dv:jumps-and-levy-models:2}
Montrer que $\phi(-i)=1$ pour la fonction caractéristique de Merton, et expliquer ce que cela signifie.
\end{exercise}

\begin{exercise}[$\star$]\label{exo:dv:jumps-and-levy-models:3}
Un modèle de Lévy a une asymétrie de $-0.4$ à un an. Calculer son asymétrie à une semaine et à quatre ans.
\end{exercise}

\begin{exercise}[$\star\star$]\label{exo:dv:jumps-and-levy-models:4}
Dériver le ratio de couverture binomial sur un événement à deux issues, et le calculer pour le straddle de biotechnologie à partir des
valeurs après l'événement~: $V_u-V_0=2.21$, $V_d-V_0=-1.48$, mouvements de $+22.1\%$ et $-14.8\%$ pour un cours au comptant de 50.
\end{exercise}

\begin{exercise}[$\star\star$]\label{exo:dv:jumps-and-levy-models:5}
Expliquer pourquoi le delta propre au modèle pour le straddle, $-0.02$, couvre plus mal que l'absence totale de couverture sur l'événement~?
\end{exercise}

\begin{exercise}[$\star\star$]\label{exo:dv:jumps-and-levy-models:6}
Expliquer pourquoi un modèle de Lévy calibré sur le smile de volatilité à un mois sous-estime le skew à un an, et ce qu'apporte le \omterm{def:dv:jumps-and-levy-models:bates}{modèle de Bates}.
\end{exercise}

\begin{exercise}[$\star\star\star$]\label{exo:dv:jumps-and-levy-models:7}
\emph{Programmation.} Relancer l'expérience de couverture quotidienne en divisant par deux la taille des sauts de Merton et en doublant leur fréquence
($\mu_J$ et $\delta$ divisés par deux, $\lambda$ doublé). Préciser comment évolue l'écart-type du P\&L couvert.
\end{exercise}

\begin{exercise}[$\star\star\star$]\label{exo:dv:jumps-and-levy-models:8}
\emph{Trouver l'erreur.} «~Le straddle à une semaine se négocie à une volatilité implicite de 160~\%. Nous l'avons vendu et couvert en delta
avec le delta de Black à 160~\%, si bien que notre seul risque est l'écart entre la volatilité réalisée et la volatilité implicite.~»
\end{exercise}

\section{Problème~: l'événement binaire}

\begin{problem}[{Problème du week-end --- ce que le smile à une semaine révèle du jeudi}]\label{pb:dv:jumps-and-levy-models:1}
Une action se traite à 50. Son smile à une semaine, coté sur 13 strikes de 38 à 62, encadre la publication des
résultats d'un essai clinique. Le modèle du desk est une diffusion plus un saut logarithmique programmé~: hausse de $a$ avec la probabilité $p$, baisse du
montant équitable dans le cas contraire. Le «~marché~» est le modèle plus riche du chapitre, dans lequel les amplitudes des issues sont elles-mêmes
incertaines.

\medskip\noindent\textbf{Partie I --- Le smile.}
\begin{enumerate}
\item Donner la volatilité implicite à la monnaie du «~marché~», ainsi que sa volatilité aux strikes 38 et 62.
\item Expliquer pourquoi le smile remonte de part et d'autre de la monnaie.
\item Donner la \omterm{def:dv:parametrising-the-surface:event}{variance d'événement} et l'écart-type du mouvement que la règle du chapitre 8 déduit de la volatilité à la
monnaie face à une diffusion de 35~\%.
\item Donner le mouvement absolu attendu que fournit la règle, et expliquer pourquoi c'est le plus fiable de ses deux nombres ici.
\item Expliquer pourquoi une seule volatilité ne permet pas de reproduire les prix des 13 strikes.
\end{enumerate}

\medskip\noindent\textbf{Partie II --- L'ajustement.}
\begin{enumerate}[resume]
\item Calibrer $(\sigma,p,a)$~: donner les paramètres et l'erreur.
\item Donner les mouvements implicites à la hausse et à la baisse.
\item Expliquer pourquoi la volatilité de diffusion ajustée est supérieure aux 35~\% réels.
\item Donner la probabilité de succès que le modèle ajusté implique, et expliquer comment un trader utiliserait ce nombre.
\item Calculer la baisse équitable pour $p=0.4$ et une hausse de $+22.1\%$.
\end{enumerate}

\medskip\noindent\textbf{Partie III --- La couverture.}
\begin{enumerate}[resume]
\item Donner la prime du straddle et l'écart-type du P\&L non couvert à travers l'événement (amplitudes connues).
\item Donner le delta de Black et le delta du modèle pour le straddle, ainsi que l'écart-type du P\&L pour chacun.
\item Donner le ratio de couverture binomial et le résidu avec des amplitudes connues.
\item Donner le meilleur ratio et le résidu lorsque les amplitudes sont incertaines, en prix et en pourcentage de la prime.
\item Indiquer quel instrument couvrirait le reste.
\end{enumerate}

\medskip\noindent\textbf{Partie IV --- Jugement.}
\begin{enumerate}[resume]
\item Expliquer pourquoi le ratio binomial n'est pas le delta de l'option.
\item Préciser comment le desk doit couvrir dans les jours qui précèdent l'événement, puis le jour de l'événement.
\item Indiquer la réserve à constituer pour un straddle vendu à travers l'événement.
\item Énoncer le \emph{résultat nommé}~: la taille et la probabilité des sauts implicites dans le smile à une semaine, ainsi que l'erreur
de la couverture en delta à travers l'événement.
\item En une phrase, dire ce qu'un modèle à sauts apporte au desk que la volatilité implicite ne peut pas fournir.
\end{enumerate}
\end{problem}

\section{Questions d'entretien}

\begin{interviewq}[$\star$ \iqroles{trader, chercheur}]\label{iq:dv:jumps-and-levy-models:1}
Pourquoi les options à courte échéance hors de la monnaie suggèrent-elles la présence de sauts~?
\end{interviewq}

\begin{interviewq}[$\star$ \iqroles{chercheur}]\label{iq:dv:jumps-and-levy-models:2}
Écrire le \omterm{def:dv:jumps-and-levy-models:merton}{modèle de diffusion à sauts de Merton} et son prix comme une série de prix de Black--Scholes.
\end{interviewq}

\begin{interviewq}[$\star\star$ \iqroles{trader, risque}]\label{iq:dv:jumps-and-levy-models:3}
Pouvez-vous couvrir un saut en delta~? Qu'est-ce qui reste, et comment le couvririez-vous~?
\end{interviewq}

\begin{interviewq}[$\star\star$ \iqroles{chercheur}]\label{iq:dv:jumps-and-levy-models:4}
Comment l'asymétrie et l'aplatissement varient-ils avec la maturité dans un modèle de Lévy, et qu'implique cela pour le smile de volatilité~?
\end{interviewq}

\begin{interviewq}[$\star\star$ \iqroles{développeur}]\label{iq:dv:jumps-and-levy-models:5}
Votre moteur de pricing de Fourier est précis pour Heston à un an et faux pour la variance gamma à une semaine. Pourquoi~?
\end{interviewq}

\begin{interviewq}[$\star\star\star$ \iqroles{trader}]\label{iq:dv:jumps-and-levy-models:6}
Une action publie ses résultats demain. Comment évaluez-vous et couvrez-vous le straddle à une semaine, et qu'est-ce qui peut mal tourner~?
\end{interviewq}
